\subsection{Ordinary trace classes with prescribed support}

Let $B$ be a connected, closed, oriented real PL four-manifold.
Keep the ordinary TOP and jointly PL family conventions, and the
subgroups $I_B,J_B$, of
Theorem~\ref{thm:four-dimensional-component-image}.
For a based ordinary loop $\gamma$ in $A_{\mathrm{TOP}}(B)$, set
\[
 T_\gamma(x,t)=(\gamma_t(x),t),\qquad
 \eta_B[\gamma]=\kappa_B(T_\gamma)
                   \in H^2(B;\mathbb Z/2).
\]
Homotopies of based loops give relative-boundary isotopies of
their traces. Additivity of the relative PL difference class,
and multiplication of loops, make $\eta_B$ a homomorphism.
It vanishes on the image of $\pi_1A_{\mathrm{PL}}(B)$, since
a jointly PL loop has a jointly PL trace. These assertions use
no assumption on the fundamental group.

\begin{theorem}\label{thm:supported-four-dimensional-traces}
Let $j:W\hookrightarrow B$ be a compact, connected, simply
connected, codimension-zero PL submanifold with collared boundary.
Allow $\partial W=\varnothing$, in which case $W=B$.
Let
\[
 j_!:H^2(W,\partial W;\mathbb Z/2)
                       \longrightarrow H^2(B;\mathbb Z/2)
\]
be extension by zero. Then
\begin{equation}\label{eq:supported-four-trace-image}
 \operatorname{im}\eta_B\supseteq\operatorname{im}j_!
 =\operatorname{PD}_B^{-1}
   \operatorname{im}\bigl(j_*:H_2(W;\mathbb Z/2)
                                  \to H_2(B;\mathbb Z/2)\bigr).
\end{equation}
Each class in this subgroup has a representing ordinary based
TOP loop which is the identity outside $\operatorname{int}W$
and on one common collar of $\partial W$ for the entire loop.
There is no assembly hypothesis on $\pi_1(B)$.

If $j_*$ on $H_2(-;\mathbb Z/2)$ is surjective, then
\begin{equation}\label{eq:supported-four-component-image}
 \operatorname{im}\eta_B=J_B=H^2(B;\mathbb Z/2),
 \qquad \operatorname{im}(\pi_0c_P)=I_B.
\end{equation}
If this cohomology group is nonzero, the automorphism comparison
is not surjective on $\pi_1$.
\end{theorem}

\begin{proof}
We retain all boundary faces in the realization argument.
Write $N=W\times I$, so that
\[
 \partial N=(W\times\partial I)\cup(\partial W\times I).
\]
Use compatible collars at the corners, or round the corners once.
The relative PL-structure convention is obtained by attaching
exterior collars and working on the resulting open
five-manifold, with the standard structure fixed near a closed
set. This is the relative formulation in
\cite[Essay~V, Section~3, equation~(3.2), p.~240]{KirbySiebenmann1977}.
For $F:N\to N$ fixing $\partial N$ pointwise, the relative
difference from the standard product PL structure $\alpha_N$ is
\begin{equation}\label{eq:relative-four-supported-class}
 \kappa_W(F)=[F_*\alpha_N]-[\alpha_N]
 \in H^3(N,\partial N;\mathbb Z/2)
       \cong H^2(W,\partial W;\mathbb Z/2).
\end{equation}
Common-collar normalization makes this construction invariant
under relative-boundary isotopy. A bottom-fixed homeomorphism
of $N$ acts trivially on its absolute homology: the inclusion
of the bottom is a homotopy equivalence and is fixed by the
homeomorphism. Poincar\'e--Lefschetz duality over $\mathbb Z/2$
then gives trivial action on the relative group in
\eqref{eq:relative-four-supported-class}. Consequently the
relative affine-lift cocycle is additive and
$\kappa_W(F^{-1})=\kappa_W(F)$.

\emph{Relative realization and ordinary loops.}
Choose one compatible smoothing of $W$, using
\cite[Essay~I, p.~39]{KirbySiebenmann1977}.
The Orson--Powell--Randal-Williams conventions allow compact
four-manifolds with boundary and fix the vertical boundary
pointwise throughout
\cite[Section~1.1, p.~1; Sections~2.1--2.2, pp.~4--5]
{OrsonPowellRandalWilliams2026}.
For a boundary-fixed $F$ their relative invariant satisfies
\begin{equation}\label{eq:supported-four-ks-comparison}
 \operatorname{KS}_{\mathrm{OPR}}(F)
             =\operatorname{PD}_W\kappa_W(F)
                         \in H_2(W;\mathbb Z/2).
\end{equation}
Here is the class comparison with the vertical faces included.
Put $Z=N\times I$ and
$C=(N\times\{0\})\cup(\partial N\times I)$.
The standard lift extends over $Z$. The changed top lift gives
a class in
$H^3(\partial Z,C;\mathbb Z/2)\cong H^3(N,\partial N;\mathbb Z/2)$.
The universal Kirby--Siebenmann square identifies its obstruction
with the connecting class in $H^4(Z,\partial Z;\mathbb Z/2)$
\cite[Essay~V, Section~5.6, pp.~252--253]{KirbySiebenmann1977}.
Contraction in the last interval shows $C\hookrightarrow Z$
is a homotopy equivalence, so this connecting map is the
interval-suspension isomorphism. Pulling back the standard
structure gives the class $\kappa_W(F^{-1})$, which equals
$\kappa_W(F)$ by the preceding trivial-action calculation.
The two suspensions followed by duality on $Z$ give precisely
the relative duality in \eqref{eq:supported-four-ks-comparison}.

Simple connectivity of $W$ makes both the assembly domain
$H_2(\pi_1W;\mathbb Z_{(2)})$ and the Bockstein target
$H_1(W;\mathbb Z_{(2)})$ zero. The realization theorem
\cite[Theorem~2.3(D), p.~6; Definition~5.8 and proof,
pp.~21--22]{OrsonPowellRandalWilliams2026}
therefore realizes every $z\in H_2(W;\mathbb Z/2)$ through
the block homotopy-fibre sequence. To obtain a literal
boundary-fixed representative, represent its element of $Q(W)$
by $F$ from the identity to a diffeomorphism $f$.
Exactness says that $f$ is smoothly pseudo-isotopic to the
identity relative to $\partial W$. Choose such a smooth
pseudo-isotopy $G$ with bottom $f$ and top identity.
Then $G'=G\circ(f^{-1}\times\mathrm{id}_I)$ has bottom
identity and top $f^{-1}$, and $F_0=G'F$ fixes every face
of $\partial N$. The OPR invariant is additive and is zero
on $G'$, so $\operatorname{KS}_{\mathrm{OPR}}(F_0)=z$.
By \eqref{eq:supported-four-ks-comparison}, every
$v\in H^2(W,\partial W;\mathbb Z/2)$ is thus
$\kappa_W(F_0)$ for some such $F_0$.

The ordinary topological pseudo-isotopy theorem for compact
simply connected four-manifolds, including those with boundary,
gives a path $H_s$ from the identity to $F_0$ which fixes
$(W\times\{0\})\cup(\partial W\times I)$ throughout
\cite[Theorem~1.5, p.~3; relative convention, p.~1]
{GabaiGayHartmanKrushkalPowell2026}.
Its top restrictions form a based ordinary loop $\gamma_s$
on $W$ fixing $\partial W$ pointwise. Set
\[
 T_s(x,t)=(\gamma_{st}(x),t),\qquad Q_s=T_s^{-1}H_s.
\]
Every $Q_s$ fixes the whole boundary of $N$;
$Q_0=\mathrm{id}$ and $Q_1=T_\gamma^{-1}F_0$.
Relative isotopy invariance, with common collars for this
compact family, gives $\kappa_W(Q_1)=0$. Additivity therefore
gives $\kappa_W(T_\gamma)=v$. This construction gives an ordinary TOP loop.

\emph{A common support collar.}
If the boundary is nonempty, choose a PL collar
$\partial W\times[0,1]\hookrightarrow W$.
For $0\le r\le1/2$ define an embedding $c_r:W\to W$
which is the identity off the collar and on it is
\[
 c_r(y,t)=(y,r+(1-r)t).
\]
Its image is $W$ with the open collar of width $r$ removed.
On this image conjugate $\gamma_s$ by $c_r$, and extend
by the identity over the removed collar; denote the result
by $h_{r,s}$. The two formulas agree on the interface because
$\gamma_s$ fixes $\partial W$. Each $h_{r,s}$ is a
homeomorphism, with inverse given by the same construction
from $\gamma_s^{-1}$, and
$h_{r,0}=h_{r,1}=\mathrm{id}$.

Continuity at $r=0$ can be checked uniformly. Give the compact
manifold a metric for which motion along a collar interval has
its stated length. Both $c_r$ and its inverse on $c_r(W)$
move points by at most $r$. If $\omega$ is a common modulus
of continuity for the compact families $\gamma_s$ and
$\gamma_s^{-1}$, then
\[
 \sup_{s,x}d(h_{r,s}(x),\gamma_s(x))\le r+\omega(r),
 \qquad
 \sup_{s,x}d(h_{r,s}^{-1}(x),\gamma_s^{-1}(x))
                                      \le r+\omega(r).
\]
On the removed collar use the fixed boundary point with the
same collar coordinate $y$ for this estimate. Thus $h$ is
a based ordinary loop homotopy. At $r=1/2$ the loop has
one common identity collar. This construction concerns
ordinary continuous TOP families.

Reparameterize the resulting loop by a PL function
$\rho:I\to I$ which is zero near $0$ and one near $1$.
The straight homotopy from $t$ to $\rho(t)$ is relative
to the endpoints. Hence the relative trace class remains
$v$, while its trace is now the identity near every face
of $\partial N$. When $\partial W$ is empty only this
time reparameterization is needed.

\emph{Extension and its class.}
Extend the loop by the identity to $B$ and call it $\Gamma$.
Gluing on the closed sets $W$ and $B\setminus\operatorname{int}W$
gives a jointly continuous family; the same construction
for inverses verifies that it is an ordinary TOP loop.
The fixed common collar makes the relative PL lifts agree
literally on the gluing region. Naturality of their difference
therefore gives
\begin{equation}\label{eq:supported-four-extension-zero}
 \eta_B[\Gamma]=\kappa_B(T_\Gamma)
                              =j_!\kappa_W(T_\gamma)=j_!v.
\end{equation}
At the cylinder level this is extension by zero from
$H^3(N,\partial N;\mathbb Z/2)$ to
$H^3(B\times I,B\times\partial I;\mathbb Z/2)$;
interval suspension identifies it with $j_!$ in degree two.

Finally,
$\operatorname{PD}_B j_!=j_*\operatorname{PD}_W$.
For example, choose a finite triangulation compatible with
$W$ and its boundary. Extend a relative cochain by zero
and cap with the fundamental chain of $B$. Simplices outside
$W$ contribute zero, and the remaining cap is the inclusion
of the cap with the relative fundamental chain of $W$.
This proves the equality in
\eqref{eq:supported-four-trace-image} and completes the
support assertion.

If $j_*$ is onto, then $\eta_B$ is onto. Every based loop
is among the paths defining $J_B$, so $J_B$ is all of
$H^2(B;\mathbb Z/2)$. The kernel statement in
Theorem~\ref{thm:four-dimensional-component-image} gives
\eqref{eq:supported-four-component-image}. Since $\eta_B$
vanishes on jointly PL loops, its nonzero target also proves
the stated failure of surjectivity on $\pi_1$.
\end{proof}

\begin{corollary}\label{cor:connected-sum-supported-traces}
Let $B=C\#D$, where $C$ is a connected, closed, simply
connected, oriented PL four-manifold and $D$ is a connected,
closed, oriented PL four-manifold with
$H_2(D;\mathbb Z/2)=0$. Then
\eqref{eq:supported-four-component-image} holds.
In particular it holds for $C\#k(S^1\times S^3)$ for every
integer $k\ge0$.
\end{corollary}

\begin{proof}
In the connected-sum decomposition take
$W=C\setminus\operatorname{int}D^4$, with boundary $S^3$.
Van Kampen applied after adding collar neighborhoods to
$C=W\cup D^4$ gives $\pi_1(W)=\pi_1(C)=0$.
Removing the open ball does not change $H_2$.
The Mayer--Vietoris sequence for the two punctured summands
has adjacent groups $H_2(S^3;\mathbb Z/2)=0$ and
$H_1(S^3;\mathbb Z/2)=0$. Hence
\[
 H_2(B;\mathbb Z/2)\cong H_2(W;\mathbb Z/2)
                                 \oplus H_2(D;\mathbb Z/2),
\]
and $j_*$ is onto. The same sequence and the product homology
of $S^1\times S^3$ prove the finite connected-sum case;
$k=0$ is the simply connected case with $W=B$.
\end{proof}

The endpoint subgroup $I_B$ consists of the components whose top
restriction is TOP-isotopic to a PL automorphism. Determining this
subgroup and the remaining concordance groups for the nonsimply
connected examples is a separate problem.
